\section{Actual trace separation and the conditional integral bound at 269}
\label{sec:cubic-model-trace-separation-bound}

The common trace classifier can now be applied to the original integral
model.  The additional step is its precise condition \texttt{M5}:
the separating trace must be an element of the actual weight-zero
algebra, with the actual integral origin.  We supply this step and
the normalized length comparison.  The assertion about the remaining
tangent obstruction below is a sufficient condition for \texttt{M4};
we do not replace \texttt{M4} by a claimed equivalent Selmer vanishing.

\paragraph{Exact domain and retained inputs.}
Put \(p=269\), \(K=\mathbb Q(\sqrt{-3})\), and let
\(\chi:G_K\to\mu_3\) be the fixed character of
\(K(\sqrt[3]{269})/K\).  The concrete arithmetic domain in this
section is
\[
 N=4,\qquad \psi=\chi_{-4},\qquad
 e\equiv267\pmod{268},\quad e\ge1,\qquad j\in\{1,2\}.
 \tag{G113.1}
\]
Retain the original coefficient DVR \(O\subset\overline{\mathbb Q}_p\),
the literal original \texttt{GD}/\texttt{Dat}, and all coefficient,
marked-element and local hypotheses used in AC-G65--AC-G68.  In
particular any original \(e_O=1\) assumption in the old endpoint
modularity theorem is retained; no larger native coefficient ring is
substituted for \(O\).  Write
\[
 E=\operatorname{Frac}O,\qquad
 e_O=\operatorname{length}_O(O/pO)>0,\qquad
 \gamma=\psi c_p^e,
 \tag{G113.2}
\]
where \(c_p(\mathrm{Frob}^{\rm geom}_w)=Q_w^{-1}\).
The residue embeddings and finite coefficient extensions used to
evaluate points are those already allowed in the actual model.

Use the actual normalized complete ring \(S_j\), its universal
semidirect representation \(r\), and the complete graded model
\[
       (A,\mathcal A,x,q_A:A\longrightarrow S_j,\mathcal Q)
 \tag{G113.3}
\]
of AC-G65-model.  Set \(I=\mathcal Q.\mathrm{ideal}\), and let
\(L_j\) be its weight-minus-one cotangent length, exactly the
quantity denoted by \(\mathcal M_j.\mathrm{len}\) there.
The actual origin is the semidirect extension of
\[
 r_*|_{G_K}=\operatorname{diag}
       (\chi^j\gamma,1,\chi^j\gamma^{-1}).
 \tag{G113.4}
\]
All structural model inputs, including \texttt{CohInp},
\texttt{LemmaP}, and the actual local component and strictness
conditions, remain as in AC-G65-model.  The two generic algebraic
inputs \texttt{LemmaM} and \texttt{LemmaM'} are retained explicitly.
They are not proved anew here.

For the arithmetic support retain one actual host and the exact
native, local and normalization hypotheses of AC-G103--AC-G110.
In particular we require both the off-triangular support assertion
(G109.4) and the native Type I support assertion (G110.2).
When these are supplied by AC-G108, its uniform endpoint condition
\((E)\), strong seed condition \((S)\), and complete native/core
premises are retained to precisely the extent used there.  The
explicit certificate AC-G103-certificate supplies the numerical gap
only on (G113.1).  A closed seed point alone is not substituted for
either support assertion.  AC-G109 then supplies \texttt{M2} over
the original \(O\), and AC-G66 supplies \texttt{M1} and
\texttt{M3} for (G113.3).

For points we retain the original-field automorphy, local
compatibility, and adjoint inputs of AC-G110 and AC-G111, and the
old actual Type II modular block and \texttt{RigiditySecondCase}
inputs of AC-G68.  For counting we retain all exact old
\texttt{Hypotheses'}, rank-two classifier and modular purity inputs
listed in AC-G112-count and AC-G112-twist.  In particular the latter
classifies \((V_f\chi^j)\oplus1\); no literal identification with
the untwisted old \texttt{Adm2B} is made.  These are the current
proved interfaces on the stated domain, not an assertion for all
Dirichlet characters or all coefficient primes.

\begin{theorem}[AC-G113-M5: the actual homogeneous trace difference]
\label{thm:G113-M5}
In the preceding domain take \(q=11\), let \(w\) be its unique
place in \(K\), put \(Q=121\), and take geometric Frobenius
\(\sigma\in G_K\) at \(w\).  There is an element
\(a_\sigma\in A\), over the original coefficient ring, such that
\[
 \begin{split}
 a_\sigma&\in\mathcal A_0,\qquad x(a_\sigma)=0,\\
 f(q_A(a_\sigma))
   &=\operatorname{tr}\rho_f(\sigma)
                -(Q^e+1+Q^{-e})
       \quad(f:S_j\to\overline{\mathbb Q}_p).
 \end{split}
 \tag{G113.5}
\]
Here \(\rho_f\) is the multiplier-one restriction to \(G_K\).
Every \(\overline{\mathbb Q}_p\)-valued point of \(A\) which
kills \(I\) and is two-sided has nonzero value on \(a_\sigma\).
Thus \(a_\sigma\) satisfies the literal source condition
\(\texttt{M5}(\mathcal A,\overline{\mathbb Q}_p,x,I,a_\sigma)\).
This conclusion does not assume \texttt{M4}.
\end{theorem}

\begin{proof}
\emph{1. The actual seed and the normalized shape alternatives.}
The rational prime 11 is inert in \(K\) and avoids \(269N3\).
The cubic character is unramified there.  A rational Frobenius
acts by inversion on its cyclic cubic quotient, since the full
Kummer splitting group is \(S_3\).  Its image outside the cubic
subgroup has square one.  Consequently Frobenius at the degree-two
place \(w\) has \(\chi(\sigma)=1\).  This is also the inert-prime
calculation of AC-K29.  The rational quadratic character \(\psi\)
has \(\psi(\sigma)=\psi(\mathrm{Frob}^{\rm geom}_{11})^2=1\).
Therefore the actual trace in (G113.4) is
\[
                   b=Q^{-e}+1+Q^e\in O.
 \tag{G113.6}
\]
It is integral over \(O\) because \(Q\) is a \(p\)-adic unit.
In fact \(b\) is a rational number already in \(\mathbb Z_p\).

We spell out the changed-tuple use of the source
\texttt{twoPt\_shape}, rather than declaring the old single-character
ring to be this new ring.  For the original marked elements \(c,g_1\),
the matrix at \(c\) is the identity, the multiplier is \(\eta_K\),
and the marked difference is normalized.  At the origin its nonzero
diagonal parameter is
\[
        \zeta_*=\chi(g_1)^j
                         (\gamma(g_1)-\gamma(g_1)^{-1})\in O^\times.
 \tag{G113.7}
\]
Every actual local point has \(\zeta_f-\zeta_*\) of \(p\)-adic
norm less than one.  Hence \(\zeta_f\ne0\).  For two local
points, \(\zeta_{f'}\ne-\zeta_f\): the opposite equality would
write the unit \(2\zeta_*\) as a sum of two elements of norm
less than one.  Here \(p\ne2\).

The complete proof of the generic matrix theorem
\texttt{shape\_cases} uses only the multiplier, normalization,
\(2\ne0\), and \(\zeta_f\ne0\).  Indeed the difference matrix
is \(\operatorname{diag}(\zeta_f,0,-\zeta_f)\); its distinct
eigenspaces make every invariant subspace a sum of
coordinate lines, and the polarization pairs the corresponding
entry vanishings.  Its four alternatives are Type I, upper,
lower, or Type II.  A two-sided actual point is neither upper
nor lower, by the actual entry grading and model signs.
It therefore has Type I or Type II.

\emph{2. Trace separation at every actual point.}
For Type I, AC-G111 and AC-G112 give the actual
original-\(K\) automorphic representation, its compact level
\(269N\), and its \texttt{Af.TrData} trace.  At \(w\) the
normalized trace has absolute value at most three under the
chosen complex coefficient isomorphism.  For Type II, AC-G68
gives \(\rho_f\simeq(V_f|_{G_K}\otimes\chi^j)\oplus1\).
The all-conjugates modular purity input of AC-G112-twist gives
the same bound, including its precise finite character label.
Thus every actual two-sided point satisfies
\[
           |\iota\operatorname{tr}\rho_f(\sigma)|\le3.
 \tag{G113.8}
\]
On the other hand \(b=Q^e+1+Q^{-e}>3\) for \(Q=121,e\ge1\).
The more general triangle bound \(Q^e-1-Q^{-e}>3\) would
also suffice.  Since every complex field isomorphism fixes the
rational number \(b\), (G113.8) proves
\(\operatorname{tr}\rho_f(\sigma)-b\ne0\).

\emph{3. Construct the element and check the literal quantifier.}
For each \(i\in\{0,1,2\}\), the actual model entry property
supplies \(E_{ii}(\sigma)\in\mathcal A_0\) whose image under
\(q_A\) is the \((i,i)\)-entry of \(r(\sigma)\).  Set
\[
        a_\sigma=E_{00}(\sigma)+E_{11}(\sigma)
                              +E_{22}(\sigma)-b\cdot1_A.
 \tag{G113.9}
\]
This is an element of the original \(A\); no scalar extension,
division by an orbit size, or change of coefficient lattice is
used.  It is homogeneous of weight zero.  The equality
\(x=x_0q_A\) with the actual basepoint proves \(x(a_\sigma)=0\),
and entrywise evaluation gives (G113.5).

Now let \(\xi:A\to\overline{\mathbb Q}_p\) kill \(I\) and
be two-sided.  The actual all-point extension property, using the
already supplied \texttt{M2}, gives \(t\in\overline{\mathbb Q}_p^\times\)
and an actual local point \(f:S_j\to\overline{\mathbb Q}_p\)
such that
\[
           \operatorname{act}(t)\xi=fq_A.
 \tag{G113.10}
\]
On a homogeneous element of weight \(d\), this action multiplies
evaluation by \(t^d\).  It preserves both sign conditions and
acts trivially on \(a_\sigma\).  The actual model signs identify
the two-sided point in (G113.10) with an actual \texttt{TwoPt}.
Consequently
\(\xi(a_\sigma)=f(q_Aa_\sigma)\ne0\), as required.  The
source \texttt{ModelSigns} proof uses the homogeneous entry lifts,
entry density, continuity of a local point, and normalization at
\(c\); none of these uses the old formula \(r_0(\eta,\Psi)\).
All four properties are fields of the actual model of AC-G65.
This verifies every clause of the literal \texttt{M5} definition.
\end{proof}

\begin{theorem}[AC-G113-bound: conditional integral length and common counting integer]
\label{thm:G113-bound}
Keep the exact domain above and assume the literal \texttt{M4}
for the model (G113.3).  Define its actual strict group
\[
 \begin{split}
 \alpha_j&=\chi^{-j}\psi^{-1}c_p^{-e}|_{G_K},\\
 \mathcal S_j&=\ker\left(H^1(G_K,(E/O)(\alpha_j))
                         \longrightarrow\mathcal L_j\right),\\
 \mathcal L_j&=\prod_{v\mid p}H^1(G_{K_v},(E/O)(\alpha_j))\\
 &\qquad\times\prod_{v\nmid p}H^1(I_v,(E/O)(\alpha_j)),\\
 s_j&=\operatorname{length}_O(\mathcal S_j)/e_O.
 \end{split}
 \tag{G113.11}
\]
Then these lengths are finite.  There is a finite set \(X_j\)
of actual Type I or cubic Type II representations, pairwise
nonisomorphic on \(G_K\), such that
\[
 s_j\ \le\ \frac{L_j}{e_O}\ \le\
 2\sum_{\rho\in X_j}
      v_p\bigl(\operatorname{tr}\rho(\sigma)-b\bigr).
 \tag{G113.12}
\]
Here \(v_p=-\log_p|\cdot|_p\), with \(v_p(p)=1\),
and \(b\) is (G113.6).  The same bound holds for any subgroup
defined by additional strict decomposition-group conditions.

Let \(C_1,A_1,c_2,A_2\) be the exact retained constants of
AC-G112-count and AC-G112-twist, and put
\[
       C=C_1+3c_2\,3^{A_2},\qquad
       A_*=A_1+A_2,\qquad Y=4(e+1).
 \tag{G113.13}
\]
Then \(\#X_j\le C Y^{A_*}\).  For any nonempty subset
\(J_0\subset\{1,2\}\) for which the displayed hypotheses,
including \texttt{M4}, hold, the common trace-data construction
gives an integer \(\mathcal N_{11,J_0}\ne0\) satisfying
\[
 \begin{split}
 \sum_{j\in J_0}s_j
       &\le2v_{269}(\mathcal N_{11,J_0}),\\
 \log|\mathcal N_{11,J_0}|
       &\le5C\,\#J_0\,Y^{A_*+1}\log11,\\
 \left(\sum_{j\in J_0}s_j\right)\log269
       &\le10C\,\#J_0\,Y^{A_*+1}\log11.
 \end{split}
 \tag{G113.14}
\]
The integer is constructed from the common finite classifier before
the coefficient prime and the sets \(X_j\) are chosen.  Its
application to integral model lengths here is only at coefficient
prime 269 in the stated domain.

A sufficient remaining arithmetic premise for \texttt{M4} is
\[
      \mathcal T_f=0
      \quad\hbox{for every actual Type II local point }f,
 \tag{G113.15}
\]
where \(\mathcal T_f\) is exactly the actual mixed tangent image
defined in (RT.6) of AC-G68, with all allowed finite coefficient
extensions.  In particular the stronger premise
\(H_f^1(K,V_f|_{G_K}\otimes\chi^j)=0\) for those points
and coefficient fields suffices.  Neither premise is proved in
this theorem.
\end{theorem}

\begin{proof}
\emph{1. Apply the generic model inequality with its exact valuation.}
The structural hypotheses of \texttt{LemmaM} are the actual fields
of (G113.3): the coefficient ring is a complete characteristic-zero
DVR, \(A\) is its stated graded algebra, \(x\) its origin, and
\(\mathcal Q\) its homogeneous equation system.  AC-G66 supplies
\texttt{M1} and \texttt{M3}, AC-G109 supplies \texttt{M2}, the
current assumption is \texttt{M4}, and AC-G113-M5 supplies the
last condition.  No old single-character \texttt{PtHyp} is being
declared for the changed tuple.

Take the source normalized valuation on the algebraic closure to be
\[
                         v_O(z)=e_O v_p(z).
 \tag{G113.16}
\]
For every nonzero \(a\in O\),
\(v_O(a)=\operatorname{length}_O(O/aO)\).  This follows by
writing \(p\) as a unit times a uniformizer to power \(e_O\);
the multiplicative and ultrametric valuation conditions are then
exactly the source \texttt{IsNVal} fields.  It is also the complete
\texttt{NormO}/\texttt{DvrFacts} normalization used in
\texttt{ModelUpper}.  In particular \(e_O\) is positive and
finite, even when the coefficient ring is ramified.

\texttt{LemmaM} now gives \(L_j<\infty\), a finite cover of
the two-sided points by torus orbits, and, for \emph{every} such
finite cover \(T\),
\[
                         L_j\le2\sum_{\xi\in T}v_O(\xi(a_\sigma)).
 \tag{G113.17}
\]
The universal quantifier on covers matters: we may replace the
first cover by an orbit-reduced cover before applying the bound.

\emph{2. Pass from orbits to pairwise nonisomorphic actual representations.}
Use (G113.10) to represent every orbit in a finite cover by an
actual local point.  Delete duplicate torus orbits, retaining a
finite family \(T_j\) which covers them all.  The homogeneous
action identity and actual model entry density show that torus
conjugacy of the resulting representations is the same orbit
relation on their \(A\)-points.  This is the generic
\texttt{ModelSigns}/\texttt{ModelUpper} cover argument.

For completeness the stronger isomorphism comparison also survives
the changed seed.  An intertwiner of the restrictions to \(G_K\)
intertwines the normalized difference matrices
\(\operatorname{diag}(\zeta_f,0,-\zeta_f)\) and
\(\operatorname{diag}(\zeta_{f'},0,-\zeta_{f'})\).
The nonzero and nonopposite conditions proved in (G113.7) force
all its off-diagonal entries to be zero.  Write its three diagonal
entries \(m_0,m_1,m_2\).  Polarization supplies both the ordinary
intertwining equations and their reverse-index versions.  If
\(m_0m_2=m_1^2\), these equations give torus conjugation with
\(t=m_0/m_1\).  Otherwise they force all middle--endpoint entries
of the representations to be zero, and a square root
\(t^2=m_0/m_2\) gives the same torus conjugation on the two
endpoint coordinates.  Such a square root exists in the algebraic
closure.  Normalization at \(c\) extends the equality from
\(G_K\) to the semidirect representation.  This is the entire
algebraic mechanism of the source \texttt{conjBy\_tor\_of\_iso};
its hypotheses do not mention a residual character formula.

Consequently distinct elements of the reduced cover give pairwise
nonisomorphic \(G_K\)-representations.  Their finite set is
\(X_j\).  On weight zero the value of \(a_\sigma\) is unchanged
through every orbit replacement.  Apply (G113.17) to this cover,
use (G113.5), and divide by the positive \(e_O\).  This gives
the second inequality in (G113.12).

\emph{3. Retain the actual integral strict group.}
AC-G65-model, through the exact AC-G52 injection and retained
\texttt{LemmaM'}, proves
\(\operatorname{length}_O\mathcal S_j\le L_j\) for the
coefficient and local conditions in (G113.11).  The p-local
condition is the full decomposition-group kernel, not merely an
inertia or ordinary quotient.  The comparison includes the
directed union of all bounded torsion levels.  Since \(L_j\) is
finite, the group has finite length.  Dividing the comparison by
the same \(e_O\) proves the first inequality in (G113.12).
Additional local kernels are subgroups and can only decrease
length.  This proof never applies the old single-character
\texttt{coefIndep} theorem to a cubic character; it stays over
the original \(O\) throughout.

\emph{4. The exact common integer and constants.}
Every element of \(X_j\) lies in the actual Type I or Type II
domain by the shape argument.  The union classifier of AC-G112
has cardinality at most \(CY^{A_*}\), with the exact factors
in (G113.13), and is injective on \(X_j\).  This proves the
claimed cardinality.

For the norm theorem use the finite index set \(J_0\), with the
root of unity \(z(j)=1\) at every label.  This is Galois stable
even when \(\#J_0=2\); distinct labels are allowed to carry
the same root.  The source product bound requires pairwise
nonisomorphic representations \emph{within} each \(X_j\),
not between different labels.  AC-G112-twist therefore supplies
one nonzero integer with
\[
 |\mathcal N_{11,J_0}|_{269}\le
       \prod_{j\in J_0}\prod_{\rho\in X_j}
          |\operatorname{tr}\rho(\sigma)-b|_{269}.
 \tag{G113.18}
\]
Taking minus logarithms reverses this norm inequality in the
correct way: the sum of the trace-difference valuations is at
most \(v_{269}(\mathcal N_{11,J_0})\).  Sum (G113.12) to
obtain the first line of (G113.14).  The exact estimate
(G112T.12), with \(q=11\) and \(\#\mathrm{Orb}=\#J_0\),
gives its second line.  Finally a nonzero rational integer has
\(v_{269}(\mathcal N)\log269\le\log|\mathcal N|\);
this proves the third line.  In particular for both cubic labels
the last constant is \(20C\), not \(10C\).

\emph{5. Which tangent premise remains.}
All actual dual-number points have the precise Bloch--Kato local
condition by AC-G66-BK, including the ramified coefficient prime.
Type I points are rigid by AC-G110 on the retained support domain.
At Type II points, the corner and middle summands vanish by the
retained old \texttt{RigiditySecondCase} and the proved Artin
calculation.  AC-G68 identifies rigidity with the vanishing of
the \emph{actual} mixed image \(\mathcal T_f\).  Thus (G113.15)
makes every actual two-sided point rigid.

Here is the precise last passage to \texttt{M4}.  Given two
derivations at a two-sided \(A\)-point killing \(I\), first
apply the torus normalization (G113.10).  On a homogeneous
weight \(d\) multiply each derivative by \(t^d\).  The product
rule and homogeneity of the generators of \(I\) show that these
are again point derivations killing \(I\).  The actual model
dual-point extension property extends them to actual dual maps
on \(S_j\).  Rigidity makes both associated cocycles coboundaries.
The source \texttt{NormCob} calculation is purely algebraic:
for a normalized representation with \(2\) and \(\zeta\)
invertible, a normalized coboundary comes from
\[
                         X\in\overline{\mathbb Q}_p
                                       \operatorname{diag}(1,0,-1).
 \tag{G113.19}
\]
It uses no old residual seed formula.  Hence the two cocycles
are linearly dependent.  The actual universal entry-density and
continuity statement makes the corresponding linear combination
of the dual derivatives zero on \(S_j\).  Restrict to \(A\)
and undo the nonzero weight scalars.  The resulting coefficients
are not both zero, and their linear combination of the original
derivations is zero.  This is exactly the definition of
\texttt{M4}, as in the full generic proof
\texttt{GradedModelDerivationConditionOfPtHyp}.

Vanishing of the whole mixed Selmer group implies (G113.15)
because every actual dual class belongs to it.  This implication
does not require, or assert, that the actual tangent image equals
the entire Selmer group.  It also does not assert that \texttt{M4}
forces either of these vanishings.
\end{proof}

\paragraph{What has been closed.}
On the concrete branch and all retained support/source hypotheses,
the actual model now has \texttt{M1}, \texttt{M2}, \texttt{M3}
and \texttt{M5}, and its trace-counting interface has the original
coefficient-ring normalization and explicit constants (G113.14).
Thus the only remaining condition among \texttt{M1}--\texttt{M5}
is \texttt{M4}; (G113.15) is a sufficient arithmetic input for it.
The strict integral comparison and the finite trace classifier
are no longer separate missing interfaces on this domain.
This does not prove mixed vanishing at 269, extend the certificate
to other \(N,\psi,e\), supply the omitted primes 2 and 3, sum
unproved model bounds over all coefficient primes, or verify the
original final numerical counting budget.  The algebraic norm
integer is coefficient independent, while the proved model
application here is expressly restricted to 269.

\paragraph{Geometric dictionary.}
The diagonal trace is an invariant function on the normalized
deformation chart.  Subtracting its value at the origin gives a
weight-zero function whose zero locus misses every two-sided
characteristic-zero point.  The retained \texttt{LemmaM} converts
that separation, together with its tangent and finiteness
hypotheses, into a cotangent-length inequality.  The cubic finite
cover and endpoint local-system twists are exactly those already
described after AC-G112-twist using the audited SGA~1 dictionary.
No new geometric object or independent vanishing theorem is being
inferred from this interpretation.
